%% ---- Supplementary Note [note:dcell-model]: DCell in TorchCell -- the model and its GO DAG ----
%% Every number here is emitted by experiments/006-kuzmin-tmi/scripts/dcell_model_go_stats.py
%% (rebuilds the filtered GO DAG as experiments/006-kuzmin-tmi/scripts/dcell.py does; writes
%% results/dcell_model/*.csv, the four panels, tab-dcell-model-go-filter.tex, and
%% tab-dcell-model-vs-paper.tex, whose published column quotes the sha256-pinned OCR text of
%% the mirror key maUsingDeepLearning2018; freezes the wandb model/params_* fields of the
%% DCell runs to dcell_wandb_model_size.csv). The figure is
%% composed by experiments/006-kuzmin-tmi/scripts/dcell_model_compose_figure.py (panel a is
%% authored there as draw.io XML). Rerun those, not this file. Detail cut from the prose on
%% 2026.09.03 (per-filter term counts, stratum scheduling, annotation-date range, per-term
%% widths, initialization) lives in tab-dcell-model-go-filter, the figure, and
%% notes/experiments.006-kuzmin-tmi.scripts.dcell_model_go_stats.md.
%% MISSING BIB KEYS (curation decision; not added here): the DCell paper (mirror key
%% maUsingDeepLearning2018: Ma et al., Nat. Methods 15, 290--298, 2018) and the Gene Ontology
%% Consortium (no key in notes/assets/bib/bib.bib yet). See REF-TO-ADD below.
\clearpage
\suppnote{note:dcell-model}{DCell in TorchCell: the ontology-structured model in the shared notation\secstatus{todo}}

%% REF-TO-ADD: maUsingDeepLearning2018 (DCell)
%% REF-TO-ADD: Gene Ontology Consortium (no key in notes/assets/bib/bib.bib yet)
DCell (Ma et al.\ 2018) is the visible-neural-network baseline of Fig.~\ref{fig:ggi}d: hidden
units grouped into subsystems, one per Gene Ontology (GO; Ashburner et al.\ 2000) term, wired
to each other as the terms are wired in the ontology. It was reimplemented in TorchCell from
the published description rather than run from the released Torch7 code, so that it reads the
same experiment records and 6,607-gene S288C reference as every other model in the
comparison; extending the published single- and double-deletion model to trigenic
interaction changes only which genes are zeroed. \suppfig{fig:dcell-model} draws the
adaptation and \supptab{tab:dcell-go-filter} the ontology it is built on; its training cost
is the subject of \suppnoteref{note:dcell-training}.

\notesec{The graph is the ontology} Where the cell encoder $F_\theta$ takes the gene--gene
relations $\{A^{(k)}\}$ only as soft priors on attention (Eq.~\ref{eq:graphreg}), DCell's
graph is the GO directed acyclic graph $G_{\mathrm{GO}}=(T,E_{\mathrm{GO}})$, edges from each
term to its parents under one synthetic super-root \texttt{GO:ROOT}, with the annotation
relation $\mathrm{genes}(t)\subseteq[N]$ of genes annotated directly to $t$; each hierarchy
edge is hard wiring, a block of a weight matrix, not a regularizer. The perturbation enters as
data: a strain with deleted set $p$ is the gene-state vector $s\in\{0,1\}^{N}$, $s_i=0$ if
$g_i\in p$, and the loader copies the reference's table of 59,986 (term, gene) rows per
strain and zeros the rows of deleted genes. DCell is therefore the rebuild route of
\suppnoteref{note:equivalence}, with no shared encoding $H$ and no operator
$\mathcal{T}_\psi$ (the $\mathcal{O}(B\,L\lvert E\rvert d)$ column of \supptab{tab:compute});
a trigenic strain zeros three rows, and that is the entire trigenic extension. Every term $t$
owns a subsystem that reads the outputs of its child subsystems $\mathrm{ch}(t)$ and the
states of its own genes,
\begin{equation}
I_t=\big[\,\Vert_{c\in\mathrm{ch}(t)}\,O_c\ \big\Vert\ s_{\mathrm{genes}(t)}\,\big],\qquad
O_t=\tanh\!\big(\mathrm{BN}(W_tI_t+b_t)\big)\in[-1,1]^{L_t},\qquad
L_t=\max\!\big(20,\lceil 0.3\,\lvert\mathrm{genes}(t)\rvert\rceil\big),
\label{eq:dcell-subsystem}
\end{equation}
with $\mathrm{BN}$ batch normalization. Two details differ from the published
$\mathrm{BatchNorm}(\mathrm{Tanh}(\mathrm{Linear}(I_t)))$: normalization precedes the
$\tanh$, and $\lvert\mathrm{genes}(t)\rvert$ counts direct annotations rather than the genes
contained by $t$ and its descendants. The root readout
$\hat y=w_r^{\top}O_{\mathrm{ROOT}}+b_r$ predicts the trigenic interaction $y_{\mathrm{int}}$
directly, with no fitness intermediate (contrast \suppnoteref{note:gi-from-fitness}); every
other subsystem carries a head $\hat y_t=w_t^{\top}O_t+b_t$, and the loss
\begin{equation}
\mathcal{L}=\mathrm{MSE}(\hat y,y)+\alpha\,\operatorname*{mean}_{t\neq r}\mathrm{MSE}(\hat y_t,y),\qquad \alpha=0.3,
\label{eq:dcell-loss}
\end{equation}
averages the auxiliary terms where the original summed them; AdamW weight decay replaces the
published $\lambda\lVert W\rVert_2$ penalty.

\notesec{Building and measuring the ontology} The published subsystems are the GO terms
left after three filters: terms with evidence code IGI (inferred from genetic interaction, a
circularity), terms containing fewer than six of the training genes, and terms redundant
with their children. TorchCell rebuilds that recipe on GO release 2025-03-16 with SGD
annotations (\supptab{tab:dcell-go-filter}): IGI-evidence annotations are removed (1,921) and
a term goes only when that empties it; a term whose gene set equals a parent's is removed;
a term containing fewer than four of the 6,607 reference genes is removed, the threshold
chosen so that the retained 2,655 subsystems approximate the published 2,526; and no date
cutoff is applied, so the ontology carries annotations that postdate the screens it is
trained on (a 2017-07-19 cutoff is implemented and would leave 1,408 subsystems, but was not
used). The result has 2,655 subsystems on 3,208 hierarchy edges in 13 strata with 59,986
annotations covering all 6,607 genes (\suppfig{fig:dcell-model}b--d). Coverage is complete
for a specific reason: 5,487 annotations with evidence code ND (no biological data) attach
genes of unknown function directly to the three namespace roots, so every gene sits in at
least two subsystems (median 8, maximum 38) and the roots are the widest subsystems
(molecular function: 2,444 direct genes, $L_t=734$). The network has 61,429 hidden units and
20,613,037 parameters, which equals the \texttt{model/params\_total} logged by every DCell
training run in both wandb projects (15 runs, one value), so the DAG measured here is the
one the baseline was trained on.

\notesec{Against the published ontology} \supptab{tab:dcell-vs-paper} sets the rebuild
beside every quantity the paper states about its ontology and model, each quoted from the
paper's text; the two are not the same graph. The paper reports 2,526 subsystems in ``12
layers'' with 97,181 neurons ``ranging from 20 to 1,075 per system''; the rebuild has 2,655
subsystems in 13 strata with 61,429 hidden units from 20 to 734. The paper reports neither
hierarchy edges, nor leaves, nor a parameter count, so the 3,208 edges, 1,527 leaves, and
20,613,037 parameters here have no published counterpart. Four causes are identifiable, one
of them measured. The GO release differs: the paper names none (it cites the 2016 Consortium
paper) and the rebuild uses the 2025-03-16 release, whose retained annotations are dated up
to 2023-10-04; under the same three filters, keeping only annotations dated on or before
2017-07-19 leaves 1,408 subsystems rather than 2,655 (\supptab{tab:dcell-go-filter}), so
release drift alone moves the count by more than the gap to 2,526. The containment filter
differs in threshold and in population: the paper drops terms ``containing fewer than six
yeast genes disrupted in the available genotypes'', the rebuild terms containing fewer than
four of all 6,607 reference genes, a threshold chosen to land near 2,526 on the newer
release. The filter semantics and order differ: IGI removal acts on annotations rather than
terms, redundancy is judged against a parent rather than the children, and the rebuild
applies redundancy before containment where the paper lists containment second. The width
rule counts direct annotations rather than contained genes, which is why the rebuild has
more subsystems but fewer hidden units: 2,521 of the 2,655 subsystems sit at the floor of
20, and \texttt{GO:ROOT}, which carries no direct annotation, is 20 wide, whereas the
published rule makes the widest subsystem 1,075 neurons, a term containing 3,581 to 3,583
genes by that rule. Whether the published ``12 layers'' counts the root is not stated. A
term-by-term comparison would need the published term list, which the paper does not
print; it was not made.

\begin{figure*}[t]
\tcfig{figures/FigS-dcell-model.pdf}{\textbf{DCell as implemented in TorchCell.} \textbf{a},~The
filtered GO DAG that structures the model, drawn whole: 2,655 subsystems as points in 13
strata (stratum $0$ is \texttt{GO:ROOT}, stratum $1$ the three namespace roots BP, MF, CC;
the deepest terms at the bottom), colored by namespace, with the 3,208 hierarchy edges
(\texttt{is\_child\_of}, child to parent) in gray; within a stratum, terms are ordered by the
mean position of their parents. The perturbation enters as data: the bottom row is the
gene-state vector $s$ over the 6,607 genes, present $=1$, and a strain with deleted set $p$
has $s_i=0$ for $g_i\in p$; the loader copies the reference's table of 59,986 (term, gene)
annotation rows per strain and zeros the rows of the deleted genes, so each strain is a
rebuilt input rather than an operator on a shared encoding. One measured strain of the
Kuzmin 2018 screen is drawn, the triple deletion YKL191W, YLR172C, YNL051W (record 4 of the
91,050-record Kuzmin 2018-only build, chosen by rule: the first record whose three
perturbations are all deletions and whose genes each carry 5 to 12 direct annotations; here
5, 5 and 9): dashed blue, the zeroed gene states entering the subsystems that annotate them
(open circles); solid red, every hierarchy edge from those subsystems up to the root, the
subsystems the zeroed rows reach. Subsystems run stratum by stratum from the deepest term to
the root. Right: the gene-state rule, the subsystem of Eq.~\ref{eq:dcell-subsystem}, the root
readout, the per-term auxiliary heads, and the loss of Eq.~\ref{eq:dcell-loss}; weights are
initialized uniformly in $[-0.001,0.001]$ as published. \textbf{b},~Subsystems per stratum of the
filtered DAG; 1,527 of the 2,655 subsystems are leaves. \textbf{c},~Genes per subsystem: direct annotations, which set
the width $L_t$, and genes contained by the term or its descendants, the quantity the
published model sized subsystems by; the median subsystem has 6 direct and 10 contained
genes. \textbf{d},~Number of subsystems each of the 6,607 genes is annotated to; 675 genes are in the
hierarchy through ND-evidence annotations to the namespace roots alone. The DAG in a and
panels b--d are measured by
\texttt{experiments/006-kuzmin-tmi/scripts/dcell\_model\_go\_stats.py}; the figure is
composed by \texttt{dcell\_model\_compose\_figure.py}.}{fig:dcell-model}
\end{figure*}

%% AUTO-GENERATED -- do not hand-edit.
%% SOURCE: experiments/006-kuzmin-tmi/scripts/dcell_model_go_stats.py
%%         reads experiments/006-kuzmin-tmi/results/dcell_model/go_filter_stages.csv
\begin{table}[!htbp]
\centering
\small
\begin{tabular}{@{}lrrrrr@{}}
\toprule
Stage & Terms & Edges & Annotations & Genes covered & Leaves \\
\midrule
SGD annotations on GO, three namespaces under \texttt{GO:ROOT} & 5,660 & 6,674 & 66,404 & 6,607 & 3,826 \\
drop IGI-evidence annotations (empty terms removed) & 5,541 & 6,540 & 64,483 & 6,607 & 3,763 \\
drop terms whose gene set equals a parent's & 5,435 & 6,433 & 64,324 & 6,607 & 3,763 \\
drop terms containing $<4$ genes (the trigenic run) & 2,655 & 3,208 & 59,986 & 6,607 & 1,527 \\
reference only: annotations dated $\le$ 2017-07-19, then the same three filters & 1,408 & 1,714 & 22,978 & 6,409 & 911 \\
\bottomrule
\end{tabular}
\caption{The GO DAG at each filtering stage of the DCell baseline. Filters are applied in the order listed (\texttt{torchcell.graph.filter\_go\_IGI}, \texttt{filter\_redundant\_terms}, \texttt{filter\_by\_contained\_genes}), on GO release 2025-03-16 with SGD gene annotations; a removed term's children are reconnected to its parents. Annotations are direct (term, gene) pairs over the 6,607-gene reference; genes covered are those with at least one such pair; leaves are terms with no child term. The last row is not part of the trigenic run.}
\label{tab:dcell-go-filter}
\end{table}


%% AUTO-GENERATED -- do not hand-edit.
%% SOURCE: experiments/006-kuzmin-tmi/scripts/dcell_model_go_stats.py
%%         published column: verbatim quotes from the OCR text of the mirror key
%%         maUsingDeepLearning2018 (paper.md, sha256 ac837bc358ea4969a72789e66e31380bfdcbd7b8aea98dec47b108ee21d2070c), listed in
%%         results/dcell_model/dcell_vs_paper.csv; TorchCell column: go_filter_stages.csv,
%%         go_terms_final.csv, go_strata.csv, go_genes_final.csv, go_edges_final.csv, dcell_model_size.csv.
%% GO release: [Introduction; reference 9] the Gene Ontology (GO), a literature-curated reference database from which we extracted 2,526 cellular subsystems (intracellular components, processes or functions)9 [...] 9. The Gene Ontology Consortium. Expansion of the Gene Ontology knowledgebase and resources. Nucleic Acids Res. 45, D331–D338 (2016).
%% Annotation source: [Online Methods, Preparation of ontologies] a biological ontology consisting of terms representing cellular subsystems, child–parent relations representing containment of one term by another, and gene-to-term annotations
%% Genes: [Online Methods, DCell architecture and training algorithm (OCR of the inline math kept as is); Preparation of ontologies] For each sample i, $X _ { i } \in \ : R ^ { M }$ denotes the genotype, represented as a binary vector of states on $M$ genes $\mathbf { \nabla } \cdot \mathbf { 1 } =$ disrupted; $0 =$ wild type) [...] 2. Terms containing fewer than six yeast genes disrupted in the available genotypes
%% Evidence filter: [Online Methods, Preparation of ontologies] 1. Terms with the evidence code ‘inferred by genetic interaction’ (IGI), to avoid potential circularity in predicting genetic interactions in the genotype–phenotype samples.
%% Redundancy filter: [Online Methods, Preparation of ontologies] 3. Terms that are redundant with respect to their children terms in the ontology.
%% Containment threshold: [Online Methods, Preparation of ontologies] 2. Terms containing fewer than six yeast genes disrupted in the available genotypes (with ‘containment’ defined as all genes annotated to that term or its descendants).
%% Filter order: [Online Methods, Preparation of ontologies] We used the following criteria to filter (remove) terms from GO: 1. Terms with the evidence code ‘inferred by genetic interaction’ (IGI), to avoid potential circularity in predicting genetic interactions in the genotype–phenotype samples. 2. Terms containing fewer than six yeast genes disrupted in the available genotypes (with ‘containment’ defined as all genes annotated to that term or its descendants). 3. Terms that are redundant with respect to their children terms in the ontology.
%% Removed-term rewiring: [Online Methods, Preparation of ontologies] When a term was removed, all children were connected directly to all parent terms to maintain the hierarchical structure.
%% Subsystems: [Online Methods, Preparation of ontologies] The remaining 2,526 terms were used to define the hierarchy of DCell subsystems.
%% Hierarchy edges: not stated in the paper
%% Depth: [Results, DCell design] The depth of both networks is 12 layers, on par with deep neural networks in other fields7.
%% Leaves: not stated in the paper
%% Width rule: [Online Methods, equation (2) (OCR of the display math kept as is)] L _ { o } ^ { ( t ) } = \operatorname* { m a x } \left( 2 0 , \left\lceil 0 . 3 * \mathrm { n u m b e r ~ o f ~ g e n e s ~ c o n t a i n e d ~ b y ~ } t \right\rceil \right)
%% Root readout: [Results, DCell design; Online Methods, equation (3)] the output layer, or root, is a single neuron representing cell phenotype [...] Linear in equation (3) denotes linear functions transforming multidimensional vector $O _ { i } ^ { ( t ) }$ into a scalar.
%% Hidden units: [Results, DCell design] The use of multiple neurons (ranging from 20 to 1,075 per system; see Online Methods) acknowledges that cellular components are often multifunctional [...] By this design, the VNN embedded in GO includes 97,181 neurons; the corresponding model for CliXO includes 22,167 neurons.
%% Parameters: not stated in the paper
%% Auxiliary loss: [Online Methods, equation (3) (OCR of the inline math kept as is)] the parameter $\alpha$ $_ { ( = 0 . 3 ) }$ balances these two contributions. $\lambda$ is an $l _ { 2 }$ norm regularization factor determined by four-fold cross-validation.
%% Training data: [Online Methods, Training genotype-phenotype data (OCR of the inline math kept as is)] Several forms of the model were employed in this study, trained on either Costanzo et al.16 (\~3 million training examples) or a more recently published update in 2016 ( ${ \sim } 8$ million training examples)15.
\begin{table}[!htbp]
\centering
\footnotesize
\begin{tabular}{@{}p{0.17\linewidth}p{0.38\linewidth}p{0.39\linewidth}@{}}
\toprule
& \textbf{Published DCell} & \textbf{TorchCell DCell} \\
\midrule
GO release & not reported; cites the GO Consortium (ref.~9, \emph{Nucleic Acids Res.} 2016) & 2025-03-16 (\texttt{go.obo} data-version) \\
Annotation source & not reported (``gene-to-term annotations'') & SGD \texttt{go\_details} over the 6,607-gene S288C reference \\
Genes & the genes disrupted in the training genotypes; count not reported & 6,607 reference genes, all covered \\
Evidence filter & terms with evidence code IGI removed & IGI annotations removed (1,921); a term goes only when emptied (119 terms) \\
Redundancy filter & terms redundant with respect to their children & terms whose gene set equals a parent's (106 terms) \\
Containment threshold & fewer than six disrupted genes, counted over the term and its descendants & fewer than four of the 6,607 reference genes, counted over the term and its descendants (2,780 terms) \\
Filter order & IGI, containment, redundancy (as listed) & IGI, redundancy, containment \\
Removed-term rewiring & children connected to all parents & children connected to all parents \\
Subsystems & 2,526 & 2,655 \\
Hierarchy edges & not reported & 3,208 \\
Depth & 12 layers & 13 strata (\texttt{GO:ROOT} plus 12) \\
Leaves & not reported & 1,527 \\
Width rule & $\max(20,\lceil 0.3\times\text{genes contained by }t\rceil)$ & $\max(20,\lceil 0.3\,\lvert\mathrm{genes}(t)\rvert\rceil)$, direct annotations \\
Root readout & one output neuron; root width not reported & one output; \texttt{GO:ROOT} has 0 direct genes, so $L_r=20$ \\
Hidden units & 97,181 (20 to 1,075 per subsystem) & 61,429 (20 to 734 per subsystem; 2,521 subsystems at the floor of 20) \\
Parameters & not reported & 20,613,037 \\
Auxiliary loss & $\alpha=0.3$, summed over $t\neq r$, plus $\lambda\lVert W\rVert_2$ ($\lambda$ by four-fold CV) & $\alpha=0.3$, averaged over $t\neq r$; AdamW weight decay \\
Training data & Costanzo 2010 ($\sim$3 M examples) or Costanzo 2016 ($\sim$8 M), single and double deletions & the trigenic dataset of the CGT experiment (\suppnoteref{note:dcell-training}) \\
\bottomrule
\end{tabular}
\caption{The published DCell ontology and model (Ma et al.\ 2018) against the TorchCell rebuild. Published values are quoted from the paper's text and Online Methods (the OCR text of the mirrored paper, sha256-pinned in the generating script); ``not reported'' marks a quantity the paper does not state. TorchCell values are measured on the frozen DAG of \supptab{tab:dcell-go-filter}; counts in parentheses are the terms or annotations each filter removed. Depth counts strata of the longest path from \texttt{GO:ROOT}; the paper's ``12 layers'' does not say whether the root is counted. Hidden units are the sum of subsystem widths $L_t$.}
\label{tab:dcell-vs-paper}
\end{table}

